\chapter{The Six Ancient Cosmologies}

This chapter covers six cosmologies. The \textbf{Greek}, the \textbf{Roman}, the
\textbf{Norse}, the \textbf{Egyptian}, the \textbf{Mesopotamian}, and the
\textbf{Chinese}. They are older and more concrete than the systems set out
earlier in this collection. They do not begin from a single abstract principle
and reason downward. They begin from a \textbf{god} or a \textbf{place} and
proceed as a narrative. The cosmos here is not derived by argument. It is
generated, through successive divine parents and children who quarrel, marry,
kill, and rule. The actors are named beings with traits and motives rather than
stages of a proof.

The same structural elements recur across all six. Each opens with a
\textbf{first emptiness} or a \textbf{first water}, a state before order with no
maker behind it. Each then performs a \textbf{separation of sky and earth}, the
splitting of an undivided whole into two halves. Each runs through a
\textbf{succession of divine generations}, an older order replaced by a younger
one, often by force. Each ends in an \textbf{ordering of the cosmos}, a settled
rule under a chief god who divides the realms and assigns the powers. The names
differ between traditions. The structure does not. The same underlying pattern
runs through all of them. The one differentiates into the many, and the many
settle into an order.
